%% Annexe B : Référentiel des outils
\chapter{Référentiel des Outils}
\label{ann:outils}

Les outils cités sont des exemples. Le cours insiste sur la méthode~: l'outil
sert la lecture, il ne la remplace pas. Tout usage se fait dans un cadre
autorisé ou en laboratoire fermé.

\begin{tableaucours}{L{4.2cm}L{5.4cm}L{2.8cm}}{%
    \tableauHeader{Catégorie d'outil} &
    \tableauHeader{Usage} &
    \tableauHeader{Chapitre, lab}
}
Collecte de journaux & Centraliser, normaliser, horodater & Ch.~\ref{chap:journaux}, SAFI-LAB-02 \\
SIEM & Corréler, alerter, cas d'usage & Ch.~\ref{chap:siem}, SAFI-LAB-03 \\
Sondes IDS/IPS & Détecter réseau et hôte & Ch.~\ref{chap:idsips} \\
Acquisition forensique & Copie bit à bit, condensat & Ch.~\ref{chap:fondements}, SAFI-LAB-05 \\
Analyse système & Artefacts, chronologie & Ch.~\ref{chap:systeme}, SAFI-LAB-06 \\
Acquisition mémoire & Capture de la mémoire vive & Ch.~\ref{chap:memoire} \\
Analyse de flux & Journaux de flux exportés & Ch.~\ref{chap:reseau}, SAFI-LAB-07 \\
Référentiel ATT\&CK & Planifier, mapper, mesurer & Ch.~\ref{chap:renseignement}, SAFI-LAB-09 \\
Tableur & Chronologie, tableau de bord, couverture & Ch.~\ref{chap:gouvernance}, SAFI-LAB-10 \\
\end{tableaucours}

\begin{boxalert}
\textbf{Rappel.} Aucun de ces outils ne se dirige vers un système réel sans
autorisation. L'entraînement se fait en laboratoire fermé.
\end{boxalert}
